%%=============================================================
%% chapter4_4_sm.tex
%% Phần 4.4.5 — State machine và luồng điều khiển
%%
%% \input vào chapter4_4.tex (hoặc trực tiếp vào main.tex sau
%% chapter4_4.tex). Các màu cppblue/pygreen/warnred/orng phải
%% được định nghĩa trước (xem chapter4_4.tex).
%%=============================================================

% ----------------------------------------------------------------
\subsection{State machine và luồng điều khiển}
\label{subsec:ai3-sm}
% ----------------------------------------------------------------

\begin{figure}[htbp]
  \centering
  \begin{tikzpicture}[
    node distance=1.0cm and 2.0cm,
    st/.style={draw, rounded corners=4pt, fill=cppblue!10, draw=cppblue!55,
               text width=2.3cm, align=center, minimum height=1.0cm, font=\small},
    ai/.style={st, fill=pygreen!10, draw=pygreen!60},
    arr/.style={->, >=Stealth, thick},
    lbl/.style={font=\scriptsize, text=gray!70, align=center},
  ]

  \node[st]  (pre)  {PREINIT\\{\scriptsize home pose}};
  \node[ai,  right=of pre]  (ini)  {INIT\\{\scriptsize AI $\to$ init\\vpDot2}};
  \node[st,  right=of ini]  (jnt)  {JOINT\\{\scriptsize vpDot2\\tracking}};
  \node[st,  right=of jnt]  (app)  {APPROACH\\{\scriptsize chờ RC5}};
  \node[st,  below=1.8cm of app]   (grp)  {GRIPPER\\{\scriptsize đóng\\gripper}};
  \node[st,  left=of grp]          (cls)  {CLASSIFIED\\{\scriptsize vận\\chuyển}};

  \draw[arr](pre)--(ini)  node[lbl,above,midway]{home\\reached};
  \draw[arr](ini)--(jnt)  node[lbl,above,midway]{vpDot2\\init ok};
  \draw[arr](jnt)--(app)  node[lbl,above,midway]{$|e|<5{\times}10^{-3}$};
  \draw[arr](app)--(grp)  node[lbl,right,midway]{RC5\\signal};
  \draw[arr](grp)--(cls)  node[lbl,below,midway]{grasped};
  \draw[arr](cls.west) to[out=180,in=270,looseness=0.8] (pre.south)
        node[lbl,left,pos=0.5]{cycle\\restart};

  \node[ai, below=1.2cm of jnt] (rec)
        {AI re-init\\{\scriptsize dot.initTracking\\(I, ai\_hint)}};

  \draw[arr, warnred!70, dashed]
        (jnt.south) -- (rec.north)
        node[midway, right, font=\scriptsize, text=warnred!70, align=left]
        {vpTracking-\\Exception};
  \draw[arr, pygreen!70, dashed]
        (rec.west) to[out=180,in=225,looseness=1.2] (jnt.south west)
        node[midway, left, font=\scriptsize, text=pygreen!70]
        {init xong};

  \draw[arr, warnred!70, dashed]
        (ini.south) -- ++(0,-0.7)
        node[below, font=\scriptsize, text=warnred]
        {fallback: click thủ công};

  \end{tikzpicture}
  \caption{State machine với vai trò rõ ràng: AI (xanh lá) cung cấp
           vị trí khởi đầu cho vpDot2 trong INIT, và re-init tracker khi
           vpDot2 báo \texttt{vpTrackingException}. vpDot2 là
           cơ chế tracking chính trong JOINT.}
  \label{fig:ai3-sm}
\end{figure}

\paragraph{PREINIT.}
Robot gửi lệnh joint position về home pose
$[0^\circ, 0^\circ, 90^\circ, 0^\circ, 90^\circ, 0^\circ]$.
Đây là tư thế chuẩn để camera có tầm nhìn tốt nhất xuống bàn làm
việc. Servo task được cấu hình với interaction matrix Desired
($\mathbf{L}_{s^*}$), pseudo-inverse, gain $\lambda=0.4$.

\paragraph{INIT.}
Chờ robot và gripper báo sẵn sàng. Khi đạt, gọi
\texttt{detectCylinderWithAI()} không có hint để lấy vị trí
cylinder, rồi khởi tạo blob tracker:
\begin{lstlisting}[language=C++]
vpImagePoint ai_hint;
if (detectCylinderWithAI(I, ai_hint)) {
    dot.initTracking(I, ai_hint);  // AI-guided init
} else {
    dot.initTracking(I);           // fallback: click thu cong
}
cog = dot.getCog();
vpFeatureBuilder::create(p, cam, dot);
\end{lstlisting}

Nếu AI thành công, tracker được khởi tạo tự động ngay tại tọa độ
pixel mà AI phát hiện. Nếu AI thất bại (timeout 3\,s hoặc không
tìm thấy cylinder), hệ thống rơi vào fallback yêu cầu người vận
hành click thủ công --- giống hệ thống ban đầu.

\paragraph{JOINT.}
vpDot2 chạy mỗi frame camera để bám theo cylinder:
\begin{lstlisting}[language=C++]
try {
    dot.track(I);           // vpDot2: nhanh, per-frame
    cog = dot.getCog();
} catch (const vpTrackingException &e) {
    // Tracking that bai --- goi AI de re-init
    vpImagePoint ai_hint;
    if (detectCylinderWithAI(I, ai_hint)) {
        dot.initTracking(I, ai_hint);
    } else {
        dot.initTracking(I);  // fallback click
    }
    cog = dot.getCog();
    vpFeatureBuilder::create(p, cam, dot);
}
vpFeatureBuilder::create(p, cam, dot);
// ... compute control law, send joint velocities
\end{lstlisting}

Mỗi vòng lặp, visual feature error được tính là:
\[
  e_x = x - x^*, \quad e_y = y - y^*
\]
trong đó $(x,y)$ là tọa độ cylinder hiện tại trong không gian ảnh
normalized (đơn vị không thứ nguyên, thu được qua phép chiếu
$x = (u - c_x)/f_x$), còn $(x^*, y^*) = (0, 0)$ là vị trí mong
muốn --- tâm ảnh. Ngưỡng hội tụ $5{\times}10^{-3}$ tương đương
khoảng $2$--$3$ pixel với thông số camera hiện tại, đủ chính xác
để gripper tiếp cận đúng vị trí cylinder.

Khi cả hai điều kiện $|e_x| < 5{\times}10^{-3}$ và
$|e_y| < 5{\times}10^{-3}$ đồng thời thỏa mãn, vòng lặp servo
dừng lại và hệ thống gửi chuỗi \texttt{"OKE\textbackslash r"} qua
cổng UART đến RC5 controller --- đây là tín hiệu thông báo
\textit{``cylinder đã được căn giữa, robot sẵn sàng để tiếp
cận''}.

\paragraph{APPROACH.}
Sau khi nhận được \texttt{"OKE\textbackslash r"}, RC5 controller
điều khiển cơ cấu hạ thấp đầu công tác xuống phía cylinder.
Trạng thái APPROACH trên C++ chỉ đơn giản là \textbf{chờ} ---
vòng lặp đọc UART và block cho đến khi RC5 gửi tín hiệu xác nhận
hoàn thành hành trình hạ (tín hiệu ``approach done''). Trong thời
gian này robot giữ nguyên joint position, camera vẫn nhìn thẳng
vào cylinder.

\paragraph{GRIPPER.}
RC5 kích hoạt đóng gripper quanh thân cylinder. Khi cảm biến
lực/vị trí xác nhận gripper đã kẹp chắc, RC5 gửi tín hiệu
``grasped'' lên C++, và hệ thống chuyển sang trạng thái
CLASSIFIED.

\paragraph{CLASSIFIED.}
Cylinder đã được gắp. Robot thực hiện hành trình vận chuyển
(joint trajectory) đến vị trí đặt xuống định trước. Sau khi đặt
xong và gripper mở ra, hệ thống quay lại PREINIT để bắt đầu
chu kỳ mới --- toàn bộ quá trình lại được thực hiện tự động mà
không cần can thiệp của người vận hành.
